\IfFileExists{stacks-project.cls}{%
\documentclass{stacks-project}
}{%
\documentclass{amsart}
}

% For dealing with references we use the comment environment
\usepackage{verbatim}
\newenvironment{reference}{\comment}{\endcomment}
%\newenvironment{reference}{}{}
\newenvironment{slogan}{\comment}{\endcomment}
\newenvironment{history}{\comment}{\endcomment}

% For commutative diagrams we use Xy-pic
\usepackage[all]{xy}

% We use 2cell for 2-commutative diagrams.
\xyoption{2cell}
\UseAllTwocells

% We use multicol for the list of chapters between chapters
\usepackage{multicol}

% This is generally recommended for better output
\usepackage{lmodern}
\usepackage[T1]{fontenc}

% For cross-file-references
\usepackage{xr-hyper}

% Package for hypertext links:
\usepackage{hyperref}

% For any local file, say "hello.tex" you want to link to please
% use \externaldocument[hello-]{hello}
\externaldocument[introduction-]{introduction}
\externaldocument[conventions-]{conventions}
\externaldocument[sets-]{sets}
\externaldocument[categories-]{categories}
\externaldocument[topology-]{topology}
\externaldocument[sheaves-]{sheaves}
\externaldocument[sites-]{sites}
\externaldocument[stacks-]{stacks}
\externaldocument[fields-]{fields}
\externaldocument[algebra-]{algebra}
\externaldocument[brauer-]{brauer}
\externaldocument[homology-]{homology}
\externaldocument[derived-]{derived}
\externaldocument[simplicial-]{simplicial}
\externaldocument[more-algebra-]{more-algebra}
\externaldocument[smoothing-]{smoothing}
\externaldocument[modules-]{modules}
\externaldocument[sites-modules-]{sites-modules}
\externaldocument[injectives-]{injectives}
\externaldocument[cohomology-]{cohomology}
\externaldocument[sites-cohomology-]{sites-cohomology}
\externaldocument[dga-]{dga}
\externaldocument[dpa-]{dpa}
\externaldocument[sdga-]{sdga}
\externaldocument[hypercovering-]{hypercovering}
\externaldocument[schemes-]{schemes}
\externaldocument[constructions-]{constructions}
\externaldocument[properties-]{properties}
\externaldocument[morphisms-]{morphisms}
\externaldocument[coherent-]{coherent}
\externaldocument[divisors-]{divisors}
\externaldocument[limits-]{limits}
\externaldocument[varieties-]{varieties}
\externaldocument[topologies-]{topologies}
\externaldocument[descent-]{descent}
\externaldocument[perfect-]{perfect}
\externaldocument[more-morphisms-]{more-morphisms}
\externaldocument[flat-]{flat}
\externaldocument[groupoids-]{groupoids}
\externaldocument[more-groupoids-]{more-groupoids}
\externaldocument[etale-]{etale}
\externaldocument[chow-]{chow}
\externaldocument[intersection-]{intersection}
\externaldocument[pic-]{pic}
\externaldocument[weil-]{weil}
\externaldocument[adequate-]{adequate}
\externaldocument[dualizing-]{dualizing}
\externaldocument[duality-]{duality}
\externaldocument[discriminant-]{discriminant}
\externaldocument[derham-]{derham}
\externaldocument[local-cohomology-]{local-cohomology}
\externaldocument[algebraization-]{algebraization}
\externaldocument[curves-]{curves}
\externaldocument[resolve-]{resolve}
\externaldocument[models-]{models}
\externaldocument[functors-]{functors}
\externaldocument[equiv-]{equiv}
\externaldocument[pione-]{pione}
\externaldocument[etale-cohomology-]{etale-cohomology}
\externaldocument[proetale-]{proetale}
\externaldocument[relative-cycles-]{relative-cycles}
\externaldocument[more-etale-]{more-etale}
\externaldocument[trace-]{trace}
\externaldocument[crystalline-]{crystalline}
\externaldocument[spaces-]{spaces}
\externaldocument[spaces-properties-]{spaces-properties}
\externaldocument[spaces-morphisms-]{spaces-morphisms}
\externaldocument[decent-spaces-]{decent-spaces}
\externaldocument[spaces-cohomology-]{spaces-cohomology}
\externaldocument[spaces-limits-]{spaces-limits}
\externaldocument[spaces-divisors-]{spaces-divisors}
\externaldocument[spaces-over-fields-]{spaces-over-fields}
\externaldocument[spaces-topologies-]{spaces-topologies}
\externaldocument[spaces-descent-]{spaces-descent}
\externaldocument[spaces-perfect-]{spaces-perfect}
\externaldocument[spaces-more-morphisms-]{spaces-more-morphisms}
\externaldocument[spaces-flat-]{spaces-flat}
\externaldocument[spaces-groupoids-]{spaces-groupoids}
\externaldocument[spaces-more-groupoids-]{spaces-more-groupoids}
\externaldocument[bootstrap-]{bootstrap}
\externaldocument[spaces-pushouts-]{spaces-pushouts}
\externaldocument[spaces-chow-]{spaces-chow}
\externaldocument[groupoids-quotients-]{groupoids-quotients}
\externaldocument[spaces-more-cohomology-]{spaces-more-cohomology}
\externaldocument[spaces-simplicial-]{spaces-simplicial}
\externaldocument[spaces-duality-]{spaces-duality}
\externaldocument[formal-spaces-]{formal-spaces}
\externaldocument[restricted-]{restricted}
\externaldocument[spaces-resolve-]{spaces-resolve}
\externaldocument[formal-defos-]{formal-defos}
\externaldocument[defos-]{defos}
\externaldocument[cotangent-]{cotangent}
\externaldocument[examples-defos-]{examples-defos}
\externaldocument[algebraic-]{algebraic}
\externaldocument[examples-stacks-]{examples-stacks}
\externaldocument[stacks-sheaves-]{stacks-sheaves}
\externaldocument[criteria-]{criteria}
\externaldocument[artin-]{artin}
\externaldocument[quot-]{quot}
\externaldocument[stacks-properties-]{stacks-properties}
\externaldocument[stacks-morphisms-]{stacks-morphisms}
\externaldocument[stacks-limits-]{stacks-limits}
\externaldocument[stacks-cohomology-]{stacks-cohomology}
\externaldocument[stacks-perfect-]{stacks-perfect}
\externaldocument[stacks-introduction-]{stacks-introduction}
\externaldocument[stacks-more-morphisms-]{stacks-more-morphisms}
\externaldocument[stacks-geometry-]{stacks-geometry}
\externaldocument[moduli-]{moduli}
\externaldocument[moduli-curves-]{moduli-curves}
\externaldocument[examples-]{examples}
\externaldocument[exercises-]{exercises}
\externaldocument[guide-]{guide}
\externaldocument[desirables-]{desirables}
\externaldocument[coding-]{coding}
\externaldocument[obsolete-]{obsolete}
\externaldocument[fdl-]{fdl}
\externaldocument[index-]{index}

% Theorem environments.
%
\theoremstyle{plain}
\newtheorem{theorem}[subsection]{Theorem}
\newtheorem{proposition}[subsection]{Proposition}
\newtheorem{lemma}[subsection]{Lemma}

\theoremstyle{definition}
\newtheorem{definition}[subsection]{Definition}
\newtheorem{example}[subsection]{Example}
\newtheorem{exercise}[subsection]{Exercise}
\newtheorem{situation}[subsection]{Situation}

\theoremstyle{remark}
\newtheorem{remark}[subsection]{Remark}
\newtheorem{remarks}[subsection]{Remarks}

\numberwithin{equation}{subsection}

% Macros
%
\def\lim{\mathop{\mathrm{lim}}\nolimits}
\def\colim{\mathop{\mathrm{colim}}\nolimits}
\def\Spec{\mathop{\mathrm{Spec}}}
\def\Hom{\mathop{\mathrm{Hom}}\nolimits}
\def\Ext{\mathop{\mathrm{Ext}}\nolimits}
\def\SheafHom{\mathop{\mathcal{H}\!\mathit{om}}\nolimits}
\def\SheafExt{\mathop{\mathcal{E}\!\mathit{xt}}\nolimits}
\def\Sch{\mathit{Sch}}
\def\Mor{\mathop{\mathrm{Mor}}\nolimits}
\def\Ob{\mathop{\mathrm{Ob}}\nolimits}
\def\Sh{\mathop{\mathit{Sh}}\nolimits}
\def\NL{\mathop{N\!L}\nolimits}
\def\CH{\mathop{\mathrm{CH}}\nolimits}
\def\proetale{{pro\text{-}\acute{e}tale}}
\def\etale{{\acute{e}tale}}
\def\QCoh{\mathit{QCoh}}
\def\Ker{\mathop{\mathrm{Ker}}}
\def\Im{\mathop{\mathrm{Im}}}
\def\Coker{\mathop{\mathrm{Coker}}}
\def\Coim{\mathop{\mathrm{Coim}}}

% Boxtimes
%
\DeclareMathSymbol{\boxtimes}{\mathbin}{AMSa}{"02}

%
% Macros for moduli stacks/spaces
%
\def\QCohstack{\mathcal{QC}\!\mathit{oh}}
\def\Cohstack{\mathcal{C}\!\mathit{oh}}
\def\Spacesstack{\mathcal{S}\!\mathit{paces}}
\def\Quotfunctor{\mathrm{Quot}}
\def\Hilbfunctor{\mathrm{Hilb}}
\def\Curvesstack{\mathcal{C}\!\mathit{urves}}
\def\Polarizedstack{\mathcal{P}\!\mathit{olarized}}
\def\Complexesstack{\mathcal{C}\!\mathit{omplexes}}
% \Pic is the operator that assigns to X its picard group, usage \Pic(X)
% \Picardstack_{X/B} denotes the Picard stack of X over B
% \Picardfunctor_{X/B} denotes the Picard functor of X over B
\def\Pic{\mathop{\mathrm{Pic}}\nolimits}
\def\Picardstack{\mathcal{P}\!\mathit{ic}}
\def\Picardfunctor{\mathrm{Pic}}
\def\Deformationcategory{\mathcal{D}\!\mathit{ef}}
\usepackage{cleveref}
\crefname{originallemma}{original lemma}{original lemmas}
\Crefname{originallemma}{Original lemma}{Original lemmas}
\makeatletter
\newlabel{sites-lemma-morphism-of-sites-covering}{{08H2}{source}{}{sourcepointer-1}{}}
\newlabel{sites-lemma-morphism-of-sites-covering@cref}{{[originallemma][0][]08H2}{[1][0][]source}{}{}{}}
\makeatother
\crefname{originaldefinition}{original definition}{original definitions}
\Crefname{originaldefinition}{Original definition}{Original definitions}
\makeatletter
\newlabel{sites-modules-definition-site-local}{{03DL}{source}{}{sourcepointer-2}{}}
\newlabel{sites-modules-definition-site-local@cref}{{[originaldefinition][0][]03DL}{[1][0][]source}{}{}{}}
\makeatother
\crefname{originallemma}{original lemma}{original lemmas}
\Crefname{originallemma}{Original lemma}{Original lemmas}
\makeatletter
\newlabel{sites-modules-lemma-local-pullback}{{03DO}{source}{}{sourcepointer-3}{}}
\newlabel{sites-modules-lemma-local-pullback@cref}{{[originallemma][0][]03DO}{[1][0][]source}{}{}{}}
\makeatother
\crefname{originalreference}{original reference}{original references}
\Crefname{originalreference}{Original reference}{Original references}
\makeatletter
\newlabel{situation-simplicial-site}{{09WE}{source}{}{sourcepointer-4}{}}
\newlabel{situation-simplicial-site@cref}{{[originalreference][0][]09WE}{[1][0][]source}{}{}{}}
\makeatother
\title{Quasi-coherent modules on the simplicial total site are exactly the cartesian ones with quasi-coherent components}
\author{The Stacks Project Authors}
\date{stacks-project:spaces-simplicial.tex:local-snapshot-20261012; GNU FDL 1.2 or later}
\begin{document}
\maketitle
\noindent\textit{Editorial note.} All delivered mathematical text is the author's own, copied from the pinned Stacks Project chapter file sources/stacks-spaces-simplicial/pinned-original-source.tex; no statement or proof step was rewritten, reordered or removed.\par\medskip
\noindent\textit{Editorial note.} The counted proof is the author's complete argument: the closing 'We omit the details' (source line 2890) covers the check that cartesianness makes the pulled-back presentation exact: pullback of modules is right exact, so the pulled-back free presentation remains a presentation of the pulled-back cokernel, and the cartesianness isomorphisms F(varphi) identify that cokernel with F\_m componentwise; the constructions on both directions are written in full, only this routine verification is unexpanded; this lemma is a stated ingredient of card L46153 of this batch.\par\medskip
\noindent\textit{Editorial note.} The label lemma-quasi-coherent-sheaf names the counted statement, and the proof environment immediately following it in the source is delivered verbatim as its proof.\par\medskip
\noindent\textit{Editorial note.} The counted statement is carried by the author's own lemma environment in the chapter source.\par\medskip
\noindent\textit{Allowed standard prerequisites (editorial).}\begin{itemize}
\item Modules on Sites, the site-local definition of quasi-coherence. Modules on Sites chapter, \textbackslash{}label\{sites-modules-definition-site-local\}. Lines 2819-2821: the sense of quasi-coherence used on both sides of the equivalence.
\item Pullbacks of quasi-coherent modules are quasi-coherent. Modules on Sites chapter, \textbackslash{}label\{sites-modules-lemma-local-pullback\}. Lines 2829-2832: the restrictions F\_n = g\_n\textasciicircum{}*F are quasi-coherent in the forward direction.
\item Coverings of C\_n versus coverings of C\_total. Same chapter, Situation \textbackslash{}ref\{situation-simplicial-site\} and the construction of the site C\_total. Lines 2842-2844: a covering of C\_total over an object of C\_n is a covering of C\_n by construction.
\item Morphism of sites: stalks and coverings (case A). Sites chapter, \textbackslash{}label\{sites-lemma-morphism-of-sites-covering\}. Line 2868: closes the cartesianness argument in case A.
\end{itemize}\medskip
\ifdefined\hypertarget\hypertarget{sourcepointer-1}{}\fi
\noindent\textit{Source reference.} Original lemma 08H2: sources/stacks-chow/tags.txt, source lines 10765--10765. Prerequisite pointer invoked at proof line(s) 2870; it names a standard fact of a sibling chapter that the counted argument uses, and every step of the counted proof built on it is retained verbatim in the delivered body.\par\medskip
\ifdefined\hypertarget\hypertarget{sourcepointer-2}{}\fi
\noindent\textit{Source reference.} Original definition 03DL: sources/stacks-chow/tags.txt, source lines 4424--4424. Prerequisite pointer invoked at proof line(s) 2822; it names a standard fact of a sibling chapter that the counted argument uses, and every step of the counted proof built on it is retained verbatim in the delivered body.\par\medskip
\ifdefined\hypertarget\hypertarget{sourcepointer-3}{}\fi
\noindent\textit{Source reference.} Original lemma 03DO: sources/stacks-chow/tags.txt, source lines 4427--4427. Prerequisite pointer invoked at proof line(s) 2830; it names a standard fact of a sibling chapter that the counted argument uses, and every step of the counted proof built on it is retained verbatim in the delivered body.\par\medskip
\ifdefined\hypertarget\hypertarget{sourcepointer-4}{}\fi
\noindent\textit{Source reference.} Original reference 09WE: sources/stacks-spaces-simplicial/pinned-original-source.tex, source lines 888--888. Prerequisite pointer invoked at proof line(s) 2818; it names a standard fact of the same chapter that the counted argument uses, and every step of the counted proof built on it is retained verbatim in the delivered body.\par\medskip
\begin{lemma}
\label{lemma-quasi-coherent-sheaf}
In Situation \ref{situation-simplicial-site}
let $\mathcal{O}$ be a sheaf of rings on $\mathcal{C}_{total}$.
Let $\mathcal{F}$ be a sheaf of $\mathcal{O}$-modules.
Then $\mathcal{F}$ is quasi-coherent in the sense of
Modules on Sites, Definition \ref{sites-modules-definition-site-local}
if and only if $\mathcal{F}$ is cartesian
and $\mathcal{F}_n$ is a quasi-coherent $\mathcal{O}_n$-module for all $n$.
\end{lemma}

\begin{proof}
Assume $\mathcal{F}$ is quasi-coherent. Since pullbacks of
quasi-coherent modules are quasi-coherent
(Modules on Sites, Lemma \ref{sites-modules-lemma-local-pullback})
we see that $\mathcal{F}_n$ is a quasi-coherent $\mathcal{O}_n$-module
for all $n$. To show that $\mathcal{F}$ is cartesian, let $U$
be an object of $\mathcal{C}_n$ for some $n$. Let us view $U$
as an object of $\mathcal{C}_{total}$. Because $\mathcal{F}$
is quasi-coherent there exists a covering $\{U_i \to U\}$
and for each $i$ a presentation
$$
\bigoplus\nolimits_{j \in J_i} \mathcal{O}_{\mathcal{C}_{total}/U_i} \to
\bigoplus\nolimits_{k \in K_i} \mathcal{O}_{\mathcal{C}_{total}/U_i} \to
\mathcal{F}|_{\mathcal{C}_{total}/U_i} \to 0
$$
Observe that $\{U_i \to U\}$ is a covering of $\mathcal{C}_n$ by
the construction of the site $\mathcal{C}_{total}$.
Next, let $V$ be an object of $\mathcal{C}_m$ for some $m$ and let
$V \to U$ be a morphism of $\mathcal{C}_{total}$ lying over
$\varphi : [n] \to [m]$. The fibre products $V_i = V \times_U U_i$
exist and we get an induced covering $\{V_i \to V\}$ in $\mathcal{C}_m$.
Restricting the presentation above to the sites
$\mathcal{C}_n/U_i$ and $\mathcal{C}_m/V_i$ we obtain
presentations
$$
\bigoplus\nolimits_{j \in J_i} \mathcal{O}_{\mathcal{C}_m/U_i} \to
\bigoplus\nolimits_{k \in K_i} \mathcal{O}_{\mathcal{C}_m/U_i} \to
\mathcal{F}_n|_{\mathcal{C}_n/U_i} \to 0
$$
and
$$
\bigoplus\nolimits_{j \in J_i} \mathcal{O}_{\mathcal{C}_m/V_i} \to
\bigoplus\nolimits_{k \in K_i} \mathcal{O}_{\mathcal{C}_m/V_i} \to
\mathcal{F}_m|_{\mathcal{C}_m/V_i} \to 0
$$
These presentations are compatible with the map
$\mathcal{F}(\varphi) : f_\varphi^*\mathcal{F}_n \to \mathcal{F}_m$
(as this map is defined using the restriction maps of $\mathcal{F}$
along morphisms of $\mathcal{C}_{total}$ lying over $\varphi$).
We conclude that $\mathcal{F}(\varphi)|_{\mathcal{C}_m/V_i}$
is an isomorphism. As $\{V_i \to V\}$ is a covering we conclude
$\mathcal{F}(\varphi)|_{\mathcal{C}_m/V}$ is an isomorphism.
Since $V$ and $U$ were arbitrary this proves that $\mathcal{F}$ is cartesian.
(In case A use Sites, Lemma \ref{sites-lemma-morphism-of-sites-covering}.)

\medskip\noindent
Conversely, assume $\mathcal{F}_n$ is quasi-coherent
for all $n$ and that $\mathcal{F}$ is cartesian.
Then for any $n$ and object $U$ of $\mathcal{C}_n$ we
can choose a covering $\{U_i \to U\}$ of $\mathcal{C}_n$
and for each $i$ a presentation
$$
\bigoplus\nolimits_{j \in J_i} \mathcal{O}_{\mathcal{C}_m/U_i} \to
\bigoplus\nolimits_{k \in K_i} \mathcal{O}_{\mathcal{C}_m/U_i} \to
\mathcal{F}_n|_{\mathcal{C}_n/U_i} \to 0
$$
Pulling back to $\mathcal{C}_{total}/U_i$ we obtain complexes
$$
\bigoplus\nolimits_{j \in J_i} \mathcal{O}_{\mathcal{C}_{total}/U_i} \to
\bigoplus\nolimits_{k \in K_i} \mathcal{O}_{\mathcal{C}_{total}/U_i} \to
\mathcal{F}|_{\mathcal{C}_{total}/U_i} \to 0
$$
of modules on $\mathcal{C}_{total}/U_i$. Then the property that
$\mathcal{F}$ is cartesian implies that this is exact.
We omit the details.
\end{proof}

\end{document}
